\subsection{Purpose and design goals}\label{subsec:knowledge-pipeline-purpose}

The knowledge-extraction pipeline converts the paragraph-level text produced by the document-extraction pipeline in Section~\ref{section:Document_Extraction} -- separated from figures, tables, and captions -- into structured biomedical claims that a downstream consumer can inspect, audit, and aggregate. 

\medskip{}

The pipeline has four design goals: 

\begin{itemize}[label={}, leftmargin=1.5em, itemsep=0.4em]

    \item \textbf{Atomic claims.} The outputs of the pipeline are subject-relation-outcome assertions, rather than paragraph summaries. Atomicity makes the downstream comparison, grouping, and contradiction-detection stages, as well as provenance attribution, more precise.
    
    \item \textbf{Source-grounded evidence.} Each finding references the evidence from which it was extracted, and these references are preserved through every stage to maintain provenance.
    
    \item \textbf{Bounded LLM stochasticity.} API calls to generative large language models are confined to the MAP stage described in Subsection~\ref{subsec:knowledge-map}; every later stage is deterministic, concentrating the pipeline's stochastic behavior into a single controlled stage.
    
    \item \textbf{Auditability, not autonomous decision-making.} The pipeline is built to let a human reviewer check and validate the extracted claims, not to make clinical decisions on its own. Final rules are ranked by an internal confidence signal, as described in Subsection~\ref{subsec:knowledge-resolve}, but these confidence scores are not clinical probabilities, and are not intended to be treated as such. 
    
\end{itemize}